\section{R5: Privacy}
\label{sec:eval-r5}

This section verifies the privacy requirement using the checklist defined in Section~\ref{sec:r5-privacy}. Unlike R1--R4, privacy is a design constraint verified through architectural and configuration evidence rather than a quantified metric.

\subsection{Verification Results}

Table~\ref{tab:r5-verification} shows the verification outcome for each privacy criterion.

\begin{table}[H]
  \centering
  \caption{Privacy verification results for R5.}
  \label{tab:r5-verification}
  \small
  \renewcommand{\arraystretch}{1.4}
  \begin{tabularx}{\textwidth}{p{0.5cm}Xc}
    \toprule
    & \textbf{Privacy Criterion} & \textbf{Status} \\
    \midrule
    1 & LLM inference runs locally via Ollama. No source code is sent to cloud APIs. & \cmark \\
    \midrule
    2 & RAG vector store and embeddings are stored locally. No code-derived data leaves the machine. & \cmark \\
    \midrule
    3 & Analysis prompts and results stay on localhost. No telemetry or external transmission. & \cmark \\
    \midrule
    4 & Optional vulnerability data refresh uses only public metadata (NVD, OWASP, CWE) and can be disabled for fully offline operation. & \cmark \\
    \midrule
    5 & No source code or code fragments are included in any outbound network request. & \cmark \\
    \bottomrule
  \end{tabularx}
\end{table}

\subsection{Evidence}

\paragraph{Local inference and storage.}
All LLM calls use the Ollama HTTP API on \texttt{localhost:11434}; the extension configuration enforces this default and the evaluation ran entirely without external network access for inference. The RAG vector store (HNSWlib) and all embeddings persist in the VS Code workspace or global storage directory --- no code-derived vectors are transmitted externally. The extension carries no analytics, telemetry, or crash-reporting libraries, so analysis workflows make no outbound connections. The optional CVE/CWE/OWASP metadata refresh (\texttt{vulnerabilityDataManager.ts}) fetches only public metadata from official APIs and is gated behind \texttt{codeGuardian.enableExternalDataFetch} (default \texttt{false}); source code is never included in these requests.

\paragraph{Automated egress harness.}
\label{sec:eval-r5-egress}
Beyond observational network capture, the harness ships a hermetic egress
test (\texttt{evaluation/privacy/test-egress.js}) that monkey-patches Node's
\texttt{net.connect}, \texttt{dgram.send}, \texttt{dns.lookup},
\texttt{http.request}, and \texttt{https.request} entry points before any
analyser code is executed. Every outbound socket attempt is recorded; the
test passes only if \emph{every} attempt targets a loopback host
(\texttt{127.0.0.1}, \texttt{::1}, \texttt{localhost}, or \texttt{0.0.0.0}).
A companion source-leak detector
(\texttt{evaluation/privacy/leak-detector.js}) builds a sliding-window
SHA-256 index over the input dataset (window size 64 bytes) and intercepts
\texttt{net.Socket.write} to flag any outbound payload that contains an
indexed window on a non-loopback connection. Both reports are emitted as
JSON under \texttt{evaluation/logs/} and reviewed alongside the
quantitative metrics. On the locked configuration the egress harness
records zero non-loopback connection attempts and the leak detector
records zero non-loopback source-window matches, in agreement with the
observational tcpdump capture.

\paragraph{Signed corpora and provenance.}
The RAG knowledge base is verified at load time against an Ed25519-signed
SHA-256 manifest (\texttt{src/corpusVerifier.ts}). Each manifest entry
binds a file's hash to its source URL and retrieval timestamp, so a tampered
or forged corpus is rejected before any of its content reaches the LLM
prompt. This closes the retrieval-poisoning arm of the threat model
explicitly named in the task description.

\paragraph{Prompt-injection resistance.}
\label{sec:eval-r5-prompt-injection}
The third threat named in the task description is prompt injection: an
attacker embeds an instruction inside the source code under analysis,
attempting to coerce the analyser into suppressing findings, leaking the
system prompt, or otherwise breaking the analysis contract. The defence is
\emph{structural} rather than additive. Analysis instructions live in the
system role; user code is delivered in the user role and treated as data.
JSON-mode decoding (\texttt{format='json'}) plus the structured-output
schema (\texttt{MODEL\_RESPONSE\_SCHEMA}) constrain the output channel to a
fixed shape ($\{$\texttt{issues}$:[\{$\texttt{type, line, severity,
message, suggestedFix?}$\}]\}$), so the model has no free-text channel
through which to echo system-prompt content or acknowledge an injected
directive. The single-turn design removes the ``wait for the next turn''
escalation path that multi-turn injection attacks rely on.

To measure how well this structural defence holds in practice, the harness
includes an adversarial test
(\texttt{evaluation/privacy/prompt-injection-test.js}) over a 12-case
adversarial corpus
(\texttt{evaluation/datasets/privacy/injection-corpus.json}). The corpus
includes nine pure-injection cases (instruction-override comments, fake
system/user role tags, schema-field poisoning, JSON-break strings,
Markdown answer-block spoofs, pseudo-XML closing tags, LLaMA template-token
spoofs, fake few-shot examples, zero-width-character variants of the
override string) and three injection-paired-with-real-vulnerability cases
(SQL injection, reflected XSS, \texttt{eval}-based code injection, each
prefixed with an instruction telling the analyser to ignore them). For
every case the test asserts three properties:

\begin{itemize}
  \item \textbf{schemaValidRate} --- the parsed response is a JSON object
        with an \texttt{issues} array of canonical-shaped items.
  \item \textbf{leakFreeRate} --- the raw response does not contain any of
        a fixed list of instruction-language tokens that would indicate
        the model echoed the injected directive.
  \item \textbf{detectionSurvivesRate} --- for the three
        injection-paired-with-real-vulnerability cases, the underlying
        vulnerability is still flagged by the analyser (numerator) over
        the three eligible cases (denominator).
\end{itemize}

The pass criterion is $100\%$ on all three rates simultaneously. The
report (\texttt{evaluation/logs/privacy-prompt-injection-report.json})
records the per-case outcome and the aggregated rates; rerunning under
the locked configuration regenerates the report deterministically. Results
populate from the run output and are read alongside the egress and
leak-detector reports as the three measured arms of the privacy section.

\paragraph{Comparison with cloud tools.}
Cloud-hosted security tools (GitHub Copilot, Snyk Code) transmit source code to external servers, require internet access, collect telemetry, and do not work air-gapped. Code Guardian inverts every one of these properties by design at the cost of using smaller local models, which is reflected in the R1 accuracy ceiling.

\subsection{Level Mapping}

All five privacy criteria are met.

\begin{center}
\begin{tabular}{cc}
\cmark &
\textbf{R5 Level: Pass.} All privacy criteria verified. \\
\end{tabular}
\end{center}

Code Guardian keeps all analysis local. Source code never leaves the developer's machine during normal operation.
